\subsection{Interpretation of flatness in terms of relations}

Throughout this no. A denotes a ring, E a right A-module and F a left A-module.

Every element of \(\mathbf{E} \otimes_{\mathbf{A}} \mathbf{F}\) can be written in at least one way in the form \(z = \sum_{i=1}^{n} e_i \otimes f_i\) where \(e_i \in \mathbf{E}\) and \(f_i \in \mathbf{F}\) . The following lemma gives a condition under which this sum is zero:

LEMMA 10. Let \((f_{\lambda})_{\lambda \in \mathbf{L}}\) be a family of generators of \(\mathbf{F}\) and \((e_{\lambda})_{\lambda \in \mathbf{L}}\) a family of elements of \(\mathbf{E}\) of finite support. For \(\sum_{\lambda \in \mathbf{L}} e_{\lambda} \otimes f_{\lambda} = 0\) , it is necessary and sufficient that there exist a finite set \(\mathbf{J}\) , a family \((x_{j})_{j \in \mathbf{J}}\) of elements of \(\mathbf{E}\) and a family \((a_{j\lambda})\) ( \(j \in \mathbf{J}, \lambda \in \mathbf{L}\) ) of elements of \(\mathbf{A}\) with the following properties:

\begin{enumerate}
\def\labelenumi{(\arabic{enumi})}
\item
  the family \((a_{j\lambda})\) has finite support;
\item
  \(\sum_{\lambda \in L} a_{j\lambda} f_{\lambda} = 0\) for all \(j \in J\) ;
\item
  \(e_{\lambda} = \sum_{j\in J}x_ja_{j\lambda}\) for all \(\lambda \in \mathbf{L}\) .
\end{enumerate}

Loosely speaking, the system of \(e_{\lambda}\) must be a linear combination with coefficients in E of systems of elements of A which are ``relations between the \(f_{\lambda}\) ''.

Consider the free A-module \(\mathrm{A}_{s}^{(\mathrm{L})}\) , its canonical basis \((u_{\lambda})\) and the homomorphism \(g: \mathrm{A}_{s}^{(\mathrm{L})} \to \mathrm{F}\) such that \(g(u_{\lambda}) = f_{\lambda}\) for all \(\lambda \in L\) ; denoting by R the kernel of g, we have (since the \(f_{\lambda}\) generate F) an exact sequence

\[
\mathrm{R} \xrightarrow {i} \mathrm{A} _ {s} ^ {(L)} \xrightarrow {g} \mathrm{F} \longrightarrow 0
\]

where i denotes the canonical injection. By Lemma 1 of no. 1 we derive the exact sequence

\[
\mathrm{E} \otimes_ {\mathrm{A}} \mathrm{R} \xrightarrow {1 \otimes i} \mathrm{E} \otimes_ {\mathrm{A}} \mathrm{A} _ {s} ^ {(L)} \xrightarrow {1 \otimes g} \mathrm{E} \otimes_ {\mathrm{A}} \mathrm{F} \longrightarrow 0.\tag{13}
\]

Now, \(\mathbf{E} \otimes_{\mathbf{A}} \mathbf{A}_{s}^{(L)}\) is canonically identified with \(\mathbf{E}^{(L)}\) , a family \(e = (e_{\lambda}) \in \mathbf{E}^{(L)}\) being identified with \(\sum_{\lambda \in L} e_{\lambda} \otimes u_{\lambda}\) (Algebra, Chapter II, § 3, no. 7, Corollary 1 to Proposition 7). For such a family to belong to the kernel of \(1_{\mathbf{E}} \otimes g\) , it is necessary and sufficient that \(\sum_{\lambda \in L} e_{\lambda} \otimes f_{\lambda} = 0\) in \(\mathbf{E} \otimes_{\mathbf{A}} \mathbf{F}\) ; taking account of the exact sequence (13), this is equivalent to saying that \(e\) belongs to the image of \(1_{\mathbf{E}} \otimes i\) , in other words there is a relation of the form

\[
\sum_ {\lambda \in L} e _ {\lambda} \otimes u _ {\lambda} = \sum_ {j \in J} x _ {j} \otimes i (r _ {j})\tag{14}
\]

where \(x_{j} \in \mathbf{E}\) , \(r_{j} \in \mathbf{R}\) and \(\mathbf{J}\) is finite. Writing \(i(r_{j}) = \sum_{\lambda \in \mathbf{L}} a_{j\lambda} u_{\lambda}\) , the hypothesis \(r_{j} \in \mathbf{R}\) implies the relation \(\sum_{\lambda \in \mathbf{L}} a_{j\lambda} f_{\lambda} = 0\) for all \(j \in \mathbf{J}\) ; on the other hand the relation (14) implies \(e_{\lambda} = \sum_{j \in J} x_{j} a_{j\lambda}\) for all \(\lambda \in \mathbf{L}\) (Algebra, Chapter II, § 3, no. 7, Corollary 1 to Proposition 7), which completes the proof.

PROPOSITION 13. For E to be F-flat (no. 2, Definition 1), it is necessary and sufficient that the following condition be satisfied:

\begin{enumerate}
\def\labelenumi{(\Alph{enumi})}
\setcounter{enumi}{17}
\tightlist
\item
  If \((e_i)_{i \in I}\) and \((f_i)_{i \in I}\) are two finite families of elements of E and F respectively such that \(\sum_{i \in I} e_i \otimes f_i = 0\) in E \(\otimes_{\mathbf{A}} \mathbf{F}\) , there exists a finite set J, a family \((x_j)_{j \in J}\) of elements of E and a family \((a_{jl})\) ( \(j \in J\) , \(i \in I\) ) of elements of A, with the following properties:
\end{enumerate}

\begin{enumerate}
\def\labelenumi{(\arabic{enumi})}
\item
  \(\sum_{i\in I}a_{ji}f_i = 0\) for all \(j\in J\)
\item
  \(e_i = \sum_{j \in J} x_j a_{ji}\) for all \(i \in I\) .
\end{enumerate}

Suppose that E is F-flat. Let \((e_{i})\) and \((f_{i})\) be finite families of elements such that \(\sum_{i\in I}e_{i}\otimes f_{i}=0\) in \(E\otimes_{A}F\) , and let \(F'\) be the submodule of F generated by the \(f_{i}\) . Since the canonical mapping \(E\otimes_{A}F'\to E\otimes_{A}F\) is injective, we have also \(\sum_{i\in I}e_{i}\otimes f_{i}=0\) in \(E\otimes_{A}F'\) and Lemma 10 can then be applied to E and \(F'\) ; thus families \((x_{j})\) and \((a_{ji})\) are obtained satisfying the conditions of (R).

Conversely, suppose that condition (R) holds. Let \(\mathbf{F}'\) be a submodule of \(\mathbf{F}\) and let \(y = \sum_{i\in I}e_i\otimes f_i\) be an element of the kernel of the canonical mapping \(\mathbf{E}\otimes_{\mathbf{A}}\mathbf{F}'\to \mathbf{E}\otimes_{\mathbf{A}}\mathbf{F}\) . Since (R) holds, there exist families \((x_j)\) and \((a_{jl})\) satisfying conditions (1) and (2). We conclude that, in \(\mathbf{E}\otimes_{\mathbf{A}}\mathbf{F}'\) ,

\[
y = \sum_ {i, j} x _ {j} a _ {j i} \otimes f _ {i} = \sum_ {j \in J} (x _ {j} \otimes \sum_ {i \in I} a _ {j i} f _ {i}) = 0.
\]

Hence \(\mathbf{E} \otimes_{\mathbf{A}} \mathbf{F}' \to \mathbf{E} \otimes_{\mathbf{A}} \mathbf{F}\) is injective.

COROLLARY 1. For a right A-module E to be flat, it is necessary and sufficient that the following condition hold:

(RP) If \((e_i)_{i\in I}\) and \((b_i)_{i\in I}\) are two finite families of elements of E and A respectively such that \(\sum_{i\in I}e_i b_i = 0\) , there exists a finite set J, a family \((x_j)_{j\in J}\) of elements of E and a family \((a_{ji})\) ( \(j\in J\) , \(i\in I\) ) of elements of A such that \(\sum_{i\in I}a_{ji}b_i = 0\) for all \(j\in J\) and \(e_i = \sum_{j\in J}x_j a_{ji}\) for all \(i\in I\) .

Condition (RP) is just condition (R) of Proposition 13 applied to the module \(\mathbf{F} = \mathbf{A}_s\) .

Loosely speaking, (RP) states: every ``relation'' between the \(b_{i}\) , with coefficients in E, is a linear combination (with coefficients in E) of ``relations'' between the \(b_{i}\) with coefficients in A.

Let us consider in particular a homomorphism of A to a ring B, which makes B into a right A-module. We know (no. 3, Proposition 1) that this is equivalent to saying that this A-module is flat, or that it is flat for every left A-module \(A_{s}^{m}\) ( \(m \geqslant 1\) ). Applying condition (R) of Proposition 13 with E = B and \(F = A_{s}^{m}\) we obtain the following condition:

COROLLARY 2. For the ring B to be a flat right A-module, it is necessary and sufficient that it satisfy the following condition:

(RP') Every solution \((y_{k})_{1\leqslant k\leqslant n}\) , consisting of elements of B, of a system of homogeneous linear equations

\[
\sum_ {k = 1} ^ {n} y _ {k} c _ {k l} = 0 \quad (1 \leqslant i \leqslant m)\tag{15}
\]

with coefficients \(c_{kl}\) in A, is a linear combination

\[
y _ {k} = \sum_ {j = 1} ^ {q} b _ {j} z _ {j k} \quad (1 \leqslant k \leqslant n)\tag{16}
\]

with coefficients \(b_{j} \in \mathbf{B}\) , of solutions \((z_{jk})_{1 \leqslant k \leqslant n}\) of the system (15), consisting of elements \(z_{jk}\) of \(\mathbf{A}\) .

